% Epistemic Verification of Physics Claims via Hyperdimensional Structural Search
% Target: Frontiers in Blockchain (DeSci section)
% Copyright (C) 2024-2026 Tristan Stoltz / Luminous Dynamics
% SPDX-License-Identifier: CC-BY-4.0

\documentclass[letterpaper,11pt]{article}
\usepackage[utf8]{inputenc}
\usepackage{times}
\usepackage{hyperref}
\usepackage{natbib}
\usepackage{booktabs}
\usepackage{amsmath,amssymb}
\usepackage{graphicx}
\usepackage[margin=1in]{geometry}
\usepackage{multicol}

\title{Epistemic Verification of Physics Claims via Hyperdimensional Structural Search}

\author{
Tristan Stoltz \\
Luminous Dynamics \\
\texttt{tristan.stoltz@evolvingresonantcocreationism.com}
}

\date{April 2026}

\begin{document}

\maketitle

\begin{abstract}
We present an epistemic verification engine for scientific claims that combines
Hyperdimensional Computing (HDC) structural search with a three-dimensional
epistemic classification framework (LEM Cube). Our system encodes 208 physics
equations across 20 domains as 16,384-dimensional hypervectors, enabling
structural analog discovery via skeleton encoding---a technique that strips
variable names to reveal pure mathematical structure. We demonstrate the system
on 10 high-profile case studies organized in 5 paired comparisons, including
the cold fusion/NIF ignition contrast, the LK-99/nickelate superconductor
divergence, and the Sch\"on/Dias fraud pattern. The system correctly classifies
verified discoveries (LIGO, Higgs) as E4/N3/M3, debunked claims (cold fusion,
EMDrive) as E0/N0/M0, and active controversies (Hubble tension) as E4/N2/M3.
We additionally report a nuclear discovery engine calibrated against 70 AME2020
reference nuclei that predicts a superheavy stability island centered at
Z=115-120, N=180, consistent with published theoretical predictions. The
complete system---including physics catalog, discovery bridge, nuclear engine,
and decentralized science portal---is open source under AGPL-3.0.
\end{abstract}

\section{Introduction}

The reproducibility crisis in science is well-documented: up to 70\% of
biomedical studies fail replication \citep{ioannidis2005}, fraudulent
publications are outpacing legitimate growth \citep{nature2025fraud}, and 2.6\%
of 2025 papers contain AI-hallucinated citations \citep{nature2026citations}.
Existing Decentralized Science (DeSci) platforms address publishing and funding
but lack computational verification of claims against established physics
\citep{desci2025frontier}.

We present three contributions:

\begin{enumerate}
\item A \textbf{physics equation catalog} of 208 landmark equations encoded as
      16,384-dimensional continuous hypervectors (ContinuousHV) using
      Hyperdimensional Computing (HDC), enabling structural analog discovery
      across 20 physics domains.

\item A \textbf{discovery bridge} that classifies physics predictions via the
      LEM Cube---a three-dimensional epistemic framework with Empirical (E0--E4),
      Normative (N0--N3), and Materiality (M0--M3) axes.

\item A \textbf{nuclear discovery engine} calibrated against 70 AME2020
      reference nuclei, predicting superheavy element stability via
      Bethe-Weizs\"acker SEMF + Woods-Saxon shell model.
\end{enumerate}

\section{Method}

\subsection{Hyperdimensional Equation Encoding}

Each equation is encoded via four parallel channels:

\begin{enumerate}
\item \textbf{Full encoding}: Recursive AST traversal with role-binding.
      Each node type (Field, Constant, DiffOperator, Sum, Product, Exponential,
      Power) gets a deterministic role vector. Children are position-bound to
      preserve ordering.

\item \textbf{Skeleton encoding}: Names stripped, leaving pure operator
      structure. This enables cross-domain analogy: the Helmholtz and
      Schr\"odinger equations have \emph{identical skeletons} ($\nabla^2\psi +
      k^2\psi = 0$) despite different physics domains.

\item \textbf{Symmetry encoding}: Lie groups (SO(n), SU(n), U(1), Lorentz,
      Poincar\'e) and discrete symmetries (C, P, T, CPT, $\mathbb{Z}_2$)
      encoded via algebra dimension scaling.

\item \textbf{Dimensional encoding}: 7 SI base dimensions as orthogonal
      basis vectors, scaled by exponent.
\end{enumerate}

Similarity is computed as a weighted combination:
\begin{equation}
S = 0.4 \cdot S_{\text{skeleton}} + 0.3 \cdot S_{\text{symmetry}} + 0.2 \cdot S_{\text{dimensional}} + 0.1 \cdot S_{\text{full}}
\end{equation}

\subsection{LEM Cube Classification}

The LEM Cube assigns three orthogonal epistemic coordinates:

\begin{itemize}
\item \textbf{E-axis (Empirical)}: E0 (unverifiable) through E4 (publicly
      reproducible), determined by simulation confidence and data availability.
\item \textbf{N-axis (Normative)}: N0 (personal) through N3 (axiomatic),
      determined by catalog similarity---high similarity to known equations
      implies higher normative authority.
\item \textbf{M-axis (Materiality)}: M0 (ephemeral) through M3 (foundational),
      determined by the physics domain.
\end{itemize}

\section{Catalog}

The catalog contains 208 equations across 20 physics domains (Table~\ref{tab:domains}).

\begin{table}[h]
\centering
\caption{Equation catalog coverage by domain.}
\label{tab:domains}
\begin{tabular}{lr|lr}
\toprule
Domain & Count & Domain & Count \\
\midrule
Classical Mechanics & 14 & Optics & 8 \\
Electromagnetism & 13 & Condensed Matter & 14 \\
General Relativity & 12 & Particle Physics & 6 \\
Quantum Mechanics & 17 & Modified Gravity & 3 \\
Quantum Field Theory & 4 & Information Theory & 14 \\
Statistical Mechanics & 13 & Biophysics & 9 \\
Fluid Dynamics & 7 & Acoustics & 3 \\
Cosmology & 6 & Plasma Physics & 11 \\
Nuclear Physics & 9 & Mathematics & 4 \\
Thermodynamics & 18 & Astrophysics & 6 \\
\bottomrule
\end{tabular}
\end{table}

\section{Case Studies}

We demonstrate the system on 10 case studies in 5 paired comparisons
(Table~\ref{tab:cases}). Each pair contrasts a verification failure with a
verification success using the same physics.

\begin{table}[h]
\centering
\caption{Case studies with LEM classifications.}
\label{tab:cases}
\begin{tabular}{lcccp{5cm}}
\toprule
Case & E & N & M & Verdict \\
\midrule
Cold Fusion (1989) & 0 & 0 & 0 & Never replicated; Gamow factor $\sim 10^{-2700}$ \\
NIF Ignition (2022--25) & 4 & 3 & 3 & Q: 1.54 $\to$ 4.13; 8 ignitions \\
\midrule
LK-99 (2023) & 0 & 0 & 0 & Cu$_2$S impurity; debunked in 3 weeks \\
Nickelates (2023--25) & 3 & 2 & 2 & Multiple groups; SLAC ambient pressure \\
\midrule
Sch\"on (2002) & 0 & 2 & 0 & Identical noise; 25+ retractions \\
Dias (2020--23) & 0 & 2 & 0 & Same pattern; 5 retractions \\
\midrule
LIGO GW150914 & 4 & 3 & 3 & SNR=24; false alarm $< 1/203{,}000$ yr \\
Higgs boson (2012) & 4 & 3 & 3 & ATLAS+CMS; 5.9$\sigma$ \\
\midrule
Hubble Tension & 4 & 2 & 3 & Both sides E4; 5$\sigma$ discrepancy \\
EMDrive & 0 & 0 & 0 & Conservation violation; thrust=0 shielded \\
\bottomrule
\end{tabular}
\end{table}

\subsection{Structural Search Results}

For the Lazar ``Gravity-A'' claim, the top-5 structural neighbors are:
Heat Equation (0.629), Yukawa Potential (0.603), Schwarzschild Metric (0.583),
de Sitter Metric (0.581), FLRW Metric (0.579). The claim sits between nuclear
force and gravitational field equations at moderate similarity ($\sim$0.6),
indicating a structural analog exists but no exact match.

For the Art's Parts metamaterial, the nearest neighbor is the Waveguide
Dispersion Relation at 0.915 similarity---the claimed physics is textbook
optics. ORNL (2022) confirmed the physical sample fails.

\section{Nuclear Discovery}

The nuclear discovery engine combines Bethe-Weizs\"acker SEMF with Woods-Saxon
shell model corrections, calibrated against 70 AME2020 reference nuclei. A
systematic sweep of Z=110-120, N=170-190 reveals a stability island centered at
Z=115, N=180 with shell corrections -17 to -23 MeV, consistent with published
predictions by M\"oller et al. and Oganessian \& Utyonkov.

\section{Discussion}

\paragraph{Limitations.} The SEMF + simplified shell model is educational, not
research-grade; calibrated predictions carry $\sigma > 5$ MeV uncertainty for
superheavy elements. The structural search uses name-hash encoding for 106 of
208 equations, limiting skeleton clustering accuracy for those entries.

\paragraph{Comparison.} No existing DeSci platform (DeSci Labs, ResearchHub,
Bio Protocol) performs computational verification of claims against a physics
equation catalog. ResearchHub's AI reviewer catches math errors; our system
identifies structural analogs across 20 domains.

\section{Conclusion}

We have demonstrated that hyperdimensional structural search, combined with
multi-dimensional epistemic classification, provides a practical framework for
automated scientific claim verification. The system correctly distinguishes
verified discoveries, debunked claims, fraudulent publications, and active
controversies. The open-source implementation at
\url{https://desci.mycelix.net} serves as both a DeSci portal and an
educational physics resource.

\bibliographystyle{plainnat}

\begin{thebibliography}{99}
\bibitem[Ioannidis(2005)]{ioannidis2005} Ioannidis, J.P.A. (2005).
Why most published research findings are false. \textit{PLoS Medicine}, 2(8), e124.

\bibitem[Wang et~al.(2021)]{ame2020} Wang, M., Huang, W.J., et~al. (2021).
The AME2020 atomic mass evaluation. \textit{Chinese Physics C}, 45(3), 030003.

\bibitem[Alcubierre(1994)]{alcubierre1994} Alcubierre, M. (1994).
The warp drive: hyper-fast travel within general relativity. \textit{Class. Quantum Grav.}, 11, L73.

\bibitem[Kanerva(2009)]{kanerva2009hd} Kanerva, P. (2009).
Hyperdimensional computing: An introduction to computing in distributed representation. \textit{Cognitive Computation}, 1(2), 139--159.

\bibitem[Nature(2025)]{nature2025fraud} Nature (2025).
A wave of retractions is shaking physics. \textit{MIT Technology Review}.

\bibitem[Nature(2026)]{nature2026citations} Nature (2026).
Hallucinated citations in scientific papers rising sharply. \textit{Nature News}.

\bibitem[DeSci(2025)]{desci2025frontier} California Management Review (2025).
Can decentralized science be the next frontier?

\bibitem[M\"oller et~al.(2016)]{moller2016} M\"oller, P., et~al. (2016).
Nuclear ground-state masses and deformations: FRDM(2012). \textit{At. Data Nucl. Data Tables}, 109-110.

\bibitem[Oganessian \& Utyonkov(2015)]{oganessian2015} Oganessian, Yu.Ts. \& Utyonkov, V.K. (2015).
Super-heavy element research. \textit{Rep. Prog. Phys.}, 78, 036301.
\end{thebibliography}

\end{document}
